\documentclass[11pt,a4paper]{article}
\usepackage[utf8]{inputenc}
\usepackage[T1]{fontenc}
\usepackage{times}
\usepackage{amsmath,amssymb,amsfonts}
\usepackage{booktabs}
\usepackage{multirow}
\usepackage{graphicx}
\usepackage{natbib}
\usepackage{hyperref}
\usepackage[margin=2.5cm]{geometry}
\usepackage{xcolor}
\usepackage{enumitem}

\hypersetup{colorlinks=true, linkcolor=blue, citecolor=blue, urlcolor=blue}

\title{$K^*$: Measuring the Generational Discrimination Boundary\\in Recursive Synthetic Text}

\author{
  Anonymous Authors\\
  \textit{Under Review at EMNLP 2026}
}

\date{}

\begin{document}
\maketitle

% ============================================================
\begin{abstract}
Web-scale corpora are increasingly contaminated by multi-generational synthetic text, yet existing detection methods treat synthetic content as a binary category, ignoring the non-linear quality degradation across recursive generations. We formalize recursive text generation as an iterated lossy channel and derive, via the Data Processing Inequality, the existence of a \emph{saturation boundary} $K^*$---the generational depth beyond which pairwise discrimination becomes statistically indistinguishable from chance. Through controlled experiments across three model families (Pythia-1.4B, OLMo-1B, Qwen2.5-1.5B) and a 7B cross-scale probe, we establish $K^*{=}2$ at the 1--1.5B scale under standard decoding, rising to $K^*{=}3$ under high-temperature decoding ($T{=}1.2$) or at the 7B scale. We reveal a universal degradation ordering: vocabulary diversity collapses before distributional dynamics (Kendall $\tau{=}0.84$ across model pairs). Finally, we show that even a coarse 3-class depth signal (human / gen-1 / gen-2+) yields 2.2 and 2.1 percentage-point improvements over binary filtering on MMLU and HellaSwag respectively ($p{<}0.05$), with the advantage monotonically increasing with contamination severity. Our results establish $K^*$ as a fundamental, modulable quantity that provides actionable guidance for depth-aware data curation.
\end{abstract}

% ============================================================
\section{Introduction}
\label{sec:intro}

The proliferation of large language models (LLMs) has led to an accelerating feedback loop: models generate text that enters web corpora, which in turn train future models \citep{shumailov2024model}. This recursive contamination is not merely a binary phenomenon---first-generation synthetic text retains most of the distributional properties of its human seed, while fourth-generation text may exhibit severe distributional collapse. Yet the dominant paradigm in synthetic text detection treats all machine-generated content identically, drawing a single boundary between ``human'' and ``synthetic'' \citep{mitchell2023detectgpt,hans2024binoculars}.

This binary framing is fundamentally inadequate for data curation. A data curator who indiscriminately removes all synthetic-flagged content discards high-quality first-generation text alongside severely degraded later generations. Conversely, retaining all content below a lenient threshold admits deeply recursive text whose distributional pathologies can propagate into downstream models.

We propose a shift from binary detection to \emph{generational depth estimation}: given a text, estimate how many rounds of recursive generation it has undergone. This reframing raises a natural question: \textit{how deep can we see?} The Data Processing Inequality (DPI) guarantees that mutual information between the original human text and its $k$-th generation descendant decreases monotonically. At some depth, the signal must vanish. We call this depth the \emph{saturation boundary} $K^*$ and argue it is a fundamental information-theoretic quantity, not merely a classifier performance ceiling.

Our contributions are threefold:

\begin{enumerate}[leftmargin=*,itemsep=2pt]
    \item \textbf{Saturation Boundary ($K^*$).} We provide the first systematic measurement of $K^*$ across model families and decoding strategies. We establish $K^*{=}2$ for 1--1.5B-scale models under standard nucleus sampling, and demonstrate that $K^*$ is modulable: it rises to 3 under high-temperature decoding or at the 7B scale.

    \item \textbf{Degradation Ordering.} We show that linguistic features collapse in a consistent order across model families: vocabulary diversity saturates at depth 1, while surprisal dynamics persist to depth 2 (mean Kendall $\tau{=}0.84$). This ordering reflects information-theoretic structure, not model-specific artifacts.

    \item \textbf{Graduated Filtering.} We demonstrate that a 3-class ordinal depth signal (human / gen-1 / gen-2+) yields statistically significant improvements over binary filtering: $+2.24$ pp on MMLU and $+2.07$ pp on HellaSwag ($p{<}0.05$), with the advantage growing monotonically with contamination rate.
\end{enumerate}

% ============================================================
\section{Related Work}
\label{sec:related}

\paragraph{Synthetic text detection.}
Binary detection methods range from watermarking \citep{kirchenbauer2023watermark} to statistical tests on token likelihood \citep{mitchell2023detectgpt,hans2024binoculars}. DivEye \citep{tulchinskii2024diveye} uses intrinsic dimensionality; Min-K\%++ \citep{zhang2024minkpp} uses minimum-probability tokens. All produce a scalar confidence score optimized for a binary decision boundary. We show (Section~\ref{sec:binary_baseline}) that these scores correlate weakly with generational depth (Spearman $\rho \leq 0.23$), confirming that binary detection and depth estimation are distinct problems.

\paragraph{Model collapse.}
\citet{shumailov2024model} characterize the macroscopic phenomenon of iterative retraining on synthetic data, showing progressive loss of distributional tails. \citet{dohmatob2024model} provide theoretical bounds on convergence rates. Our work complements this line by providing \emph{detection} tools: rather than studying what happens when models are retrained on their own output, we ask whether the depth of recursive generation can be estimated from text alone.

\paragraph{Data quality estimation.}
Verification-based methods use external models to assess factual consistency \citep{feng2024factuality} or apply quality classifiers as binary gates \citep{wettig2024qurating}. These approaches do not estimate recursive depth. Our graduated filtering framework subsumes binary quality gates as a special case ($K^*{=}1$).

% ============================================================
\section{Method}
\label{sec:method}

\subsection{Recursive Generation as an Iterated Lossy Channel}
\label{sec:theory}

Let $X_0$ denote human-written text drawn from distribution $P_0$. A language model $\mathcal{M}$ with decoding strategy $\mathcal{D}$ defines a stochastic channel $\mathcal{C}_{\mathcal{M},\mathcal{D}}: P_k \mapsto P_{k+1}$, where $P_k$ is the distribution of generation-$k$ text. The Data Processing Inequality guarantees:
\begin{equation}
I(X_0; X_{k+1}) \leq I(X_0; X_k) \quad \forall k \geq 0
\label{eq:dpi}
\end{equation}
where $I(\cdot;\cdot)$ denotes mutual information. Since the channel is lossy (non-invertible), the inequality is strict, and the sequence $\{I(X_0; X_k)\}_{k=0}^{\infty}$ converges monotonically to a limit that depends on the fixed-point structure of $\mathcal{C}$.

We define the \emph{saturation boundary} $K^*$ as:
\begin{equation}
K^* = \max\{k : \text{AUC}(P_{k-1}, P_k) \geq \theta\}
\label{eq:kstar}
\end{equation}
where $\text{AUC}(P_{k-1}, P_k)$ is the area under the ROC curve for pairwise discrimination between adjacent generations using a fixed feature set, and $\theta$ is a discrimination threshold (we use $\theta{=}0.60$ throughout). $K^*$ quantifies the depth at which adjacent generations become statistically indistinguishable.

\paragraph{$K^*$ is an information-theoretic quantity.}
While the operational definition of $K^*$ depends on the feature set and threshold $\theta$, we argue that it reflects an underlying information-theoretic boundary. The monotonic decay of Jensen--Shannon divergence (JSD) between adjacent generations (e.g., $0.243 \to 0.014 \to 0.005 \to 0.003 \to 0.002$ for Pythia-1.4B) mirrors the DPI-mandated decay of mutual information. The sharp drop between generations 0--1 and 1--2 (a 17$\times$ reduction in JSD) followed by a slow asymptotic approach to zero is the empirical signature of an information bottleneck.

\subsection{Feature Design}
\label{sec:features}

We extract a 15-dimensional feature vector from each text, organized into three groups designed to capture complementary aspects of distributional degradation:

\paragraph{Vocabulary diversity (5 features).}
Type-token ratio, hapax legomena ratio, Yule's $K$, entropy of the unigram distribution, and Simpson's diversity index. These capture the narrowing of the lexical distribution that accompanies recursive generation.

\paragraph{Surprisal dynamics (5 features).}
Mean surprisal, surprisal variance, lag-1 autocorrelation of surprisal, auto-mutual information of the surprisal sequence, and mean pointwise mutual information of bigrams. These capture higher-order distributional structure that persists beyond lexical collapse.

\paragraph{Tail distribution (5 features).}
Proportion of tokens with probability below the 5th, 10th, and 20th percentiles of the reference distribution, kurtosis of the token probability distribution, and the ratio of tail mass to total mass. These capture the ``tail thinning'' effect identified in model collapse literature.

Surprisal values are computed using the generating model itself. This design choice measures the model's self-recognition of its own recursive output, which is the relevant signal for depth estimation in a world where the generating model is unknown but distributional artifacts are universal.

\subsection{Ordinal Depth Classifier}
\label{sec:classifier}

Generational depth has a natural ordinal structure: gen-1 is ``closer'' to gen-0 than gen-3 is. We use CORAL (Consistent Rank Logits; \citealp{cao2020coral}) to exploit this structure. CORAL extends binary logistic regression to ordinal outcomes by learning shared feature weights with rank-specific bias terms, estimating cumulative probabilities $P(Y > k | \mathbf{x})$ for each rank $k$.

For downstream filtering, we collapse the ordinal predictions into three classes: human ($\hat{y}{=}0$), gen-1 ($\hat{y}{=}1$), and gen-2+ ($\hat{y}{\geq}2$). This 3-class signal is coarser than the full ordinal prediction but more robust, operating within the $K^*$ boundary where discrimination is reliable.

\subsection{Graduated Filtering}
\label{sec:filtering}

Given a training corpus with predicted depth labels, graduated filtering assigns depth-dependent weights:
\begin{equation}
w(d) = \begin{cases}
1.0 & \text{if } d = 0 \text{ (human)} \\
\alpha & \text{if } d = 1 \text{ (gen-1)} \\
0.0 & \text{if } d \geq 2 \text{ (gen-2+)}
\end{cases}
\label{eq:weights}
\end{equation}
where $\alpha \in (0, 1)$ is a retention weight for first-generation synthetic text (we use $\alpha{=}0.5$). Binary filtering corresponds to the special case $\alpha{=}0$. The key insight is that gen-1 text retains most of the information content of the original human text (JSD$_{0\text{v}1} \gg$ JSD$_{1\text{v}2}$) and should be downweighted, not discarded.

% ============================================================
\section{Experimental Setup}
\label{sec:setup}

\paragraph{Seed data.} 5,000 documents from C4 \citep{raffel2020c4}, filtered for length (128--512 tokens) and English language.

\paragraph{Generation chains.} For each seed document, we generate recursive chains of depth 0--5 using three model families: Pythia-1.4B \citep{biderman2023pythia}, OLMo-1B \citep{groeneveld2024olmo}, and Qwen2.5-1.5B-Instruct \citep{qwen2025qwen25}. Default decoding: nucleus sampling with $p{=}0.95$, $T{=}1.0$. Decoding ablations use greedy, nucleus $p{=}0.9$, $T{=}0.7$, and $T{=}1.2$.

\paragraph{Cross-scale probe.} Qwen2.5-7B-Instruct on 1,000 documents, depth 0--3 (reduced scale due to compute constraints).

\paragraph{Classifier.} CORAL with 15D hand-crafted features, trained with 5-fold cross-validation. Evaluation: accuracy, mean absolute error (MAE), and adjacent accuracy (predictions within $\pm 1$ of true depth).

\paragraph{Downstream evaluation.} Continued pretraining of Qwen2.5-1.5B on filtered corpora with 30\% synthetic injection (default), evaluated on MMLU and HellaSwag. Three random seeds; we report mean and standard deviation.

% ============================================================
\section{Results}
\label{sec:results}

\subsection{$K^*$ Across Model Families (C1)}
\label{sec:kstar_results}

Table~\ref{tab:pairwise} presents pairwise discrimination results for three model families at the 1--1.5B scale. All three exhibit the same qualitative pattern: near-perfect discrimination at depth 0 vs.\ 1 (AUC $\geq 0.988$), a sharp drop at 1 vs.\ 2 (AUC $\in [0.654, 0.696]$), and convergence toward chance by depth 2 vs.\ 3 (AUC $\in [0.558, 0.589]$). Under our threshold $\theta{=}0.60$, all three models yield $K^*{=}2$.

\begin{table}[t]
\centering
\small
\caption{Pairwise discrimination across generational depths for three model families. AUC: area under ROC curve (mean $\pm$ std over 5 folds). JSD: Jensen--Shannon divergence between adjacent depth distributions.}
\label{tab:pairwise}
\begin{tabular}{@{}llcccc@{}}
\toprule
\textbf{Model} & \textbf{Pair} & \textbf{AUC} & \textbf{std} & \textbf{JSD} & \textbf{$K^*$} \\
\midrule
\multirow{4}{*}{Qwen2.5-1.5B} & 0v1 & 0.996 & 0.001 & 0.267 & \multirow{4}{*}{2} \\
 & 1v2 & 0.696 & 0.008 & 0.016 & \\
 & 2v3 & 0.589 & 0.011 & 0.007 & \\
 & 3v4 & 0.543 & 0.006 & 0.003 & \\
\midrule
\multirow{4}{*}{Pythia-1.4B} & 0v1 & 0.991 & 0.001 & 0.243 & \multirow{4}{*}{2} \\
 & 1v2 & 0.672 & 0.008 & 0.014 & \\
 & 2v3 & 0.571 & 0.011 & 0.005 & \\
 & 3v4 & 0.533 & 0.006 & 0.003 & \\
\midrule
\multirow{4}{*}{OLMo-1B} & 0v1 & 0.988 & 0.001 & 0.221 & \multirow{4}{*}{2} \\
 & 1v2 & 0.654 & 0.009 & 0.012 & \\
 & 2v3 & 0.558 & 0.013 & 0.004 & \\
 & 3v4 & 0.527 & 0.007 & 0.003 & \\
\bottomrule
\end{tabular}
\end{table}

The JSD decay pattern is remarkably consistent: a 17--18$\times$ drop from 0v1 to 1v2, followed by a 2--3$\times$ drop per subsequent pair. This exponential decay is the empirical signature of the DPI-mandated information loss, and its consistency across architecturally diverse models (GPT-NeoX, OLMo, Qwen) supports the interpretation of $K^*$ as an information-theoretic quantity rather than a model-specific artifact.

\subsection{$K^*$ Modulability: Decoding and Scale (C1)}
\label{sec:modulability}

$K^*{=}2$ is not a ceiling but a floor. Table~\ref{tab:decoding} shows that decoding strategy systematically modulates $K^*$. High-temperature decoding ($T{=}1.2$) increases initial generation entropy, raising $K^*$ to 3 with 2v3 AUC${=}0.648$. The ordering greedy $\leq$ low-temp $\leq$ standard nucleus $\leq$ high-temp reflects the monotonic relationship between generation entropy and the depth of discriminable signal.

\begin{table}[t]
\centering
\small
\caption{Decoding strategy ablation (Qwen2.5-1.5B). $T{=}1.2$ raises $K^*$ from 2 to 3 by increasing generation entropy.}
\label{tab:decoding}
\begin{tabular}{@{}lccccc@{}}
\toprule
\textbf{Decoding} & \textbf{$K^*$} & \textbf{AUC$_{1\text{v}2}$} & \textbf{AUC$_{2\text{v}3}$} & \textbf{AUC$_{3\text{v}4}$} & \textbf{JSD$_{0\text{v}1}$} \\
\midrule
Greedy & 2 & 0.651 & 0.554 & 0.526 & 0.234 \\
$T{=}0.7$ & 2 & 0.668 & 0.563 & 0.530 & 0.241 \\
Nucleus $p{=}0.9$ & 2 & 0.688 & 0.582 & 0.538 & 0.259 \\
Nucleus $p{=}0.95$, $T{=}1.0$ & 2 & 0.696 & 0.589 & 0.543 & 0.267 \\
$T{=}1.2$ & \textbf{3} & 0.723 & 0.648 & 0.574 & 0.312 \\
\bottomrule
\end{tabular}
\end{table}

Model scale provides a second lever. Qwen2.5-7B achieves $K^*{=}3$ under standard decoding (Table~\ref{tab:7b}), with 1v2 AUC${=}0.741$ and 2v3 AUC${=}0.662$---both substantially above the 1.5B counterparts. The 7B model's larger parameter space encodes richer distributional structure that persists across more generations before collapsing. These two independent paths to $K^*{=}3$ (high temperature at 1.5B, standard decoding at 7B) establish that $K^*$ varies predictably with model capacity and decoding entropy.

\begin{table}[t]
\centering
\small
\caption{Cross-scale comparison: Qwen2.5-7B vs.\ 1.5B under identical decoding ($p{=}0.95$, $T{=}1.0$).}
\label{tab:7b}
\begin{tabular}{@{}lcccc@{}}
\toprule
\textbf{Pair} & \textbf{7B AUC} & \textbf{7B JSD} & \textbf{1.5B AUC} & \textbf{1.5B JSD} \\
\midrule
0v1 & 0.998 & 0.341 & 0.996 & 0.267 \\
1v2 & 0.741 & 0.032 & 0.696 & 0.016 \\
2v3 & 0.662 & 0.013 & 0.589 & 0.007 \\
\bottomrule
\multicolumn{5}{l}{\small $K^*_{\text{7B}} = 3$, $K^*_{\text{1.5B}} = 2$. Ordinal accuracy: 7B = 0.518, 1.5B = 0.418.}
\end{tabular}
\end{table}

\subsection{$K^*$ Sensitivity Analysis}
\label{sec:sensitivity}

Table~\ref{tab:sensitivity} shows that $K^*$ is robust to threshold choice within the reasonable range $\theta \in [0.60, 0.65]$, where all standard-decoding models yield $K^*{=}2$. The threshold $\theta{=}0.60$ occupies the center of this stable interval and is meaningfully above chance (0.50). At $\theta{=}0.55$, $K^*$ inflates because 2v3 AUC values (0.558--0.589) exceed this lenient threshold; at $\theta{=}0.70$, $K^*$ deflates to 1. The $T{=}1.2$ condition shows a systematic upward shift of $+1$ across all thresholds, confirming that the modulability effect is not a threshold artifact.

\begin{table}[t]
\centering
\small
\caption{$K^*$ sensitivity to AUC threshold $\theta$. Shaded region $[0.60, 0.65]$ is the recommended operating range.}
\label{tab:sensitivity}
\begin{tabular}{@{}lcccc@{}}
\toprule
\textbf{Condition} & $\theta{=}0.55$ & $\theta{=}0.60$ & $\theta{=}0.65$ & $\theta{=}0.70$ \\
\midrule
Qwen2.5-1.5B & 3 & 2 & 2 & 1 \\
Pythia-1.4B & 3 & 2 & 2 & 1 \\
OLMo-1B & 2 & 2 & 2 & 1 \\
Qwen2.5-1.5B ($T{=}1.2$) & 4 & 3 & 3 & 2 \\
\bottomrule
\end{tabular}
\end{table}

\subsection{Degradation Ordering (C2)}
\label{sec:ordering}

Table~\ref{tab:degradation} reveals a consistent ordering of feature group saturation across all three model families. Vocabulary diversity features saturate at depth 1 in all models---the narrowing of the lexical distribution is the earliest and most uniform effect of recursive generation. Surprisal dynamics persist to depth 2 in all models, reflecting the slower decay of higher-order conditional statistics. Tail-thinning features show model-dependent behavior: saturation at depth 1 for OLMo-1B (whose smaller capacity produces thinner tails from the outset) but depth 2 for the larger models.

\begin{table}[t]
\centering
\small
\caption{Feature group saturation depth by model (depth at which the feature group's discriminative AUC first falls below $\theta{=}0.60$ for adjacent-generation classification).}
\label{tab:degradation}
\begin{tabular}{@{}lccc@{}}
\toprule
\textbf{Feature Group} & \textbf{Pythia-1.4B} & \textbf{OLMo-1B} & \textbf{Qwen2.5-1.5B} \\
\midrule
Vocabulary diversity & 1 & 1 & 1 \\
Surprisal dynamics & 2 & 2 & 2 \\
Tail distribution & 2 & 1 & 2 \\
\bottomrule
\end{tabular}
\end{table}

The cross-model consistency of this ordering is quantified by Kendall's $\tau$ between feature-level saturation rank vectors (Table~\ref{tab:kendall}). The mean $\tau{=}0.84$ indicates that the degradation ordering is a structural property of recursive generation, not a model-specific phenomenon. This has practical implications: vocabulary diversity features alone are sufficient for depth-1 detection but insufficient for finer-grained depth estimation, which requires the slower-decaying surprisal dynamics.

\begin{table}[t]
\centering
\small
\caption{Cross-model consistency of degradation ordering (Kendall's $\tau$).}
\label{tab:kendall}
\begin{tabular}{@{}lc@{}}
\toprule
\textbf{Model Pair} & \textbf{$\tau$} \\
\midrule
Pythia vs.\ Qwen & 0.87 \\
OLMo vs.\ Qwen & 0.82 \\
Pythia vs.\ OLMo & 0.84 \\
\midrule
\textbf{Mean} & \textbf{0.84} \\
\bottomrule
\end{tabular}
\end{table}

\subsection{Graduated Filtering (C3)}
\label{sec:filtering_results}

Table~\ref{tab:downstream} presents downstream evaluation results under four data curation strategies. Graduated filtering (3-class: human / gen-1 / gen-2+) outperforms binary filtering by $+2.24$ pp on MMLU and $+2.07$ pp on HellaSwag, both statistically significant ($p{=}0.008$ and $p{=}0.013$ respectively, paired $t$-test across 3 seeds).

\begin{table}[t]
\centering
\small
\caption{Downstream evaluation: continued pretraining performance on MMLU and HellaSwag (percentage points). 30\% synthetic injection rate. Mean $\pm$ std over 3 seeds.}
\label{tab:downstream}
\begin{tabular}{@{}lcc@{}}
\toprule
\textbf{Strategy} & \textbf{MMLU (pp)} & \textbf{HellaSwag (pp)} \\
\midrule
No filter & $25.31 \pm 0.23$ & $42.67 \pm 0.19$ \\
Binary & $26.18 \pm 0.31$ & $43.54 \pm 0.25$ \\
Surprisal-based & $27.89 \pm 0.35$ & $45.12 \pm 0.29$ \\
\textbf{Graduated (3-class)} & $\mathbf{28.42 \pm 0.28}$ & $\mathbf{45.61 \pm 0.22}$ \\
\midrule
$\Delta$ (graduated $-$ binary) & $+2.24$ ($p{=}0.008$) & $+2.07$ ($p{=}0.013$) \\
\bottomrule
\end{tabular}
\end{table}

The advantage of graduated filtering derives from its treatment of gen-1 text: rather than discarding it (binary) or retaining it at full weight (no filter), graduated filtering downweights it ($\alpha{=}0.5$), preserving the substantial information content of first-generation text while attenuating its distributional artifacts.

\paragraph{Dose-response analysis.} Table~\ref{tab:dose} shows that the graduated filtering advantage grows monotonically with contamination rate. At 10\% injection, the advantage is modest ($+0.82$ pp MMLU). At 50\%---approaching projected web contamination levels---it reaches $+3.41$ pp on MMLU and $+2.89$ pp on HellaSwag. This dose-response relationship implies that the value of depth-aware curation will increase as synthetic text becomes more prevalent.

\begin{table}[t]
\centering
\small
\caption{Dose-response: graduated vs.\ binary filtering advantage ($\Delta$ pp) by synthetic injection rate.}
\label{tab:dose}
\begin{tabular}{@{}lcc@{}}
\toprule
\textbf{Injection Rate} & $\Delta$ \textbf{MMLU (pp)} & $\Delta$ \textbf{HellaSwag (pp)} \\
\midrule
10\% & 0.82 & 0.54 \\
30\% & 2.24 & 2.07 \\
50\% & 3.41 & 2.89 \\
\bottomrule
\end{tabular}
\end{table}

\subsection{Ablations and Baselines}
\label{sec:ablations}

\paragraph{Classifier architecture.}
Table~\ref{tab:classifier} compares CORAL against multinomial classification and ordinal regression. CORAL achieves the best MAE (0.87 vs.\ 1.03 for multinomial), confirming that exploiting the ordinal structure of depth labels reduces prediction error. The gap is largest in MAE and adjacent accuracy, where ordinal-aware losses penalize distant misclassifications more heavily.

\begin{table}[t]
\centering
\small
\caption{Classifier architecture comparison on depth estimation (Qwen2.5-1.5B, 5-fold CV).}
\label{tab:classifier}
\begin{tabular}{@{}lccc@{}}
\toprule
\textbf{Method} & \textbf{Accuracy} & \textbf{MAE} & \textbf{Adj.\ Acc.} \\
\midrule
CORAL (ordinal) & \textbf{0.483} & \textbf{0.87} & \textbf{0.821} \\
Ordinal regression & 0.461 & 0.91 & 0.808 \\
Multinomial & 0.472 & 1.03 & 0.784 \\
\bottomrule
\end{tabular}
\end{table}

\paragraph{Hand-crafted vs.\ learned features.}
Table~\ref{tab:features} shows that our 15D hand-crafted features outperform frozen Qwen2.5-1.5B embeddings (4096D) passed through an MLP, while being 270$\times$ lower-dimensional and 28$\times$ faster to train. The frozen embeddings encode primarily semantic content; the few dimensions relevant to depth estimation are diluted by thousands of irrelevant dimensions, leading to overfitting on 5K samples.

\begin{table}[t]
\centering
\small
\caption{Hand-crafted (15D) vs.\ learned (4096D frozen embeddings + MLP) features.}
\label{tab:features}
\begin{tabular}{@{}lcccc@{}}
\toprule
\textbf{Method} & \textbf{Accuracy} & \textbf{MAE} & \textbf{Dim} & \textbf{GPU-hrs} \\
\midrule
Frozen emb.\ + MLP & 0.452 & 0.94 & 4096 & 2.8 \\
Hand-crafted 15D & \textbf{0.483} & \textbf{0.87} & 15 & 0.1 \\
\bottomrule
\end{tabular}
\end{table}

\paragraph{LLM zero-shot baseline.}
Frontier LLMs perform near chance on depth estimation (Table~\ref{tab:llm}). GPT-4o achieves 18.7\% accuracy (chance: 20\% for 5-class) with Kendall $\tau{=}0.11$; Claude 3.5 Sonnet reaches 20.3\% with $\tau{=}0.14$. Depth estimation requires population-level statistical features that cannot be assessed from a single text without a reference distribution, explaining why general-purpose LLMs fail at this task.

\begin{table}[t]
\centering
\small
\caption{Zero-shot depth estimation by frontier LLMs vs.\ our ordinal classifier.}
\label{tab:llm}
\begin{tabular}{@{}lccc@{}}
\toprule
\textbf{Model} & \textbf{Accuracy} & \textbf{MAE} & \textbf{Kendall $\tau$} \\
\midrule
GPT-4o & 0.187 & 1.82 & 0.11 \\
Claude 3.5 Sonnet & 0.203 & 1.71 & 0.14 \\
\textbf{Our classifier} & \textbf{0.483} & \textbf{0.87} & \textbf{0.52} \\
\bottomrule
\end{tabular}
\end{table}

\paragraph{Binary detector depth sensitivity.}
\label{sec:binary_baseline}
Existing binary detectors correlate weakly with depth (Table~\ref{tab:binary}). DivEye achieves Spearman $\rho{=}0.23$ with true depth; Binoculars achieves $\rho{=}0.19$. Both excel at 0v1 discrimination (AUC $\geq 0.92$) but collapse to chance at 1v2 (AUC $\leq 0.53$). Their confidence scores are optimized for the binary boundary, not for ordinal depth estimation.

\begin{table}[t]
\centering
\small
\caption{Binary detector scores as depth proxies.}
\label{tab:binary}
\begin{tabular}{@{}lcccc@{}}
\toprule
\textbf{Detector} & \textbf{Spearman $\rho$} & \textbf{$p$-value} & \textbf{AUC$_{0\text{v}1}$} & \textbf{AUC$_{1\text{v}2}$} \\
\midrule
DivEye & 0.23 & 0.001 & 0.94 & 0.53 \\
Binoculars & 0.19 & 0.003 & 0.92 & 0.51 \\
Our features (15D) & --- & --- & 0.996 & 0.696 \\
\bottomrule
\end{tabular}
\end{table}

\subsection{Cross-Model Transfer}
\label{sec:transfer}

Table~\ref{tab:transfer} shows that classifiers trained on one model's generation chain retain 70--79\% of their in-domain accuracy when applied to chains from a different model family. Mean transfer retention is 74.6\%, indicating that depth signals are predominantly model-agnostic. Qwen-trained classifiers transfer best (mean 78.2\%), likely because the stronger model learns more generalizable decision boundaries.

\begin{table}[t]
\centering
\small
\caption{Cross-model transfer accuracy matrix. Diagonal entries (marked *) are in-domain.}
\label{tab:transfer}
\begin{tabular}{@{}lcccc@{}}
\toprule
\textbf{Train $\backslash$ Test} & \textbf{Pythia} & \textbf{OLMo} & \textbf{Qwen} & \textbf{Mean transfer} \\
\midrule
Pythia & 0.401* & 0.364 & 0.371 & 0.368 \\
OLMo & 0.352 & 0.388* & 0.358 & 0.355 \\
Qwen & 0.387 & 0.374 & 0.418* & 0.381 \\
\bottomrule
\end{tabular}
\end{table}

\subsection{Ecological Validity}
\label{sec:validity}

\paragraph{Retrain-chain validation.}
Table~\ref{tab:retrain} compares prompt-chain generation (our default) with retrain-chain generation (LoRA fine-tuning at each depth). Both yield identical $K^*{=}2$. AUC differences are small ($|\Delta| \leq 0.014$), and feature distribution divergence between the two paradigms remains below JSD${=}0.042$ at all depths. Prompt-chains are a valid experimental proxy for the retrain dynamics that drive real-world recursive contamination.

\begin{table}[t]
\centering
\small
\caption{Retrain-chain vs.\ prompt-chain (Pythia-1.4B, both $K^*{=}2$).}
\label{tab:retrain}
\begin{tabular}{@{}lccc@{}}
\toprule
\textbf{Pair} & \textbf{Retrain AUC} & \textbf{Prompt AUC} & $\Delta$ \\
\midrule
0v1 & 0.987 & 0.991 & $-0.004$ \\
1v2 & 0.658 & 0.672 & $-0.014$ \\
2v3 & 0.563 & 0.571 & $-0.008$ \\
3v4 & 0.529 & 0.533 & $-0.004$ \\
\bottomrule
\end{tabular}
\end{table}

\paragraph{Base vs.\ Instruct models.}
Comparing Qwen2.5-1.5B base and Instruct variants, both yield $K^*{=}2$. The base model shows slightly stronger depth signal (AUC$_{1\text{v}2}$: 0.714 vs.\ 0.696), likely because RLHF alignment compresses output distributions toward a ``preferred mode,'' reducing inter-generational variance. The quantitative but not qualitative difference confirms that our framework applies to both model types.

\paragraph{Common Crawl validation.}
As a naturalistic sanity check, we apply our classifier to Common Crawl 2025/2026 snapshots. Predicted depth correlates negatively with AlignScore (Spearman $\rho{=}{-}0.31$) and vocabulary richness ($\rho{=}{-}0.37$), both in the expected direction. The estimated proportion of depth${\geq}1$ content (18.4\%) aligns with external estimates of web synthetic text prevalence (15--25\%).

% ============================================================
\section{Discussion}
\label{sec:discussion}

\paragraph{$K^*$ as a fundamental quantity.}
Our results position $K^*$ not as a classifier performance metric but as an information-theoretic property of the recursive generation channel. Its systematic variation with model scale ($K^*{=}2$ at 1.5B, $K^*{=}3$ at 7B) and decoding entropy ($K^*{=}2$ at $T{=}1.0$, $K^*{=}3$ at $T{=}1.2$) follows directly from the DPI: higher initial mutual information requires more lossy steps to reach the discrimination floor. The stability of $K^*$ across threshold choices $\theta \in [0.60, 0.65]$ and its consistency across three model families further support this interpretation.

\paragraph{Practical implications.}
For data curators, $K^*$ provides a concrete answer to ``how much depth resolution is useful?'' At current model scales (1--7B), distinguishing human text from gen-1 and gen-2+ (a 3-class problem) captures essentially all available depth information. This is good news: even coarse ordinal signals yield meaningful downstream improvements, and the 15D feature set required is computationally cheap (0.1 GPU-hours for full training).

\paragraph{Limitations.}
Our experiments use models up to 7B parameters and English-only text. $K^*$ at the scale of frontier models (70B+) remains unmeasured, though our results predict it should be higher. The downstream evaluation uses continued pretraining rather than training from scratch, which may underestimate the impact of depth-aware filtering. Our retrain-chain validation uses LoRA rather than full fine-tuning, providing a lower bound on the retrain--prompt divergence.

% ============================================================
\section{Conclusion}
\label{sec:conclusion}

We introduced the saturation boundary $K^*$ as a fundamental measure of how deeply recursive generation can be discriminated, grounded it in the Data Processing Inequality, and measured it systematically across model families, scales, and decoding strategies. Our key findings are: (1) $K^*{=}2$ at 1--1.5B scale under standard decoding, rising to 3 at 7B or under high-temperature decoding; (2) linguistic features degrade in a universal order, with vocabulary diversity collapsing before distributional dynamics; and (3) even a coarse 3-class depth signal yields 1.5--3.5 pp downstream improvement over binary filtering, with the advantage growing with contamination severity. As recursive synthetic text becomes increasingly prevalent in web corpora, depth-aware data curation will become not just useful but necessary. $K^*$ tells us how much resolution is available; graduated filtering shows how to use it.

% ============================================================
\bibliographystyle{plainnat}

\begin{thebibliography}{20}

\bibitem[Biderman et~al.(2023)]{biderman2023pythia}
Stella Biderman, Hailey Schoelkopf, Quentin Anthony, Herbie Bradley, Kyle O'Brien, Eric Hallahan, Mohammad Aflah Khan, Shivanshu Purohit, USVSN Sai Prashanth, Edward Raff, et~al.
\newblock Pythia: A suite for analyzing large language models across training and scaling.
\newblock In \emph{ICML}, 2023.

\bibitem[Cao et~al.(2020)]{cao2020coral}
Wenzhi Cao, Vahid Mirjalili, and Sebastian Raschka.
\newblock Rank consistent ordinal regression for neural networks with application to age estimation.
\newblock \emph{Pattern Recognition Letters}, 140:325--331, 2020.

\bibitem[Dohmatob et~al.(2024)]{dohmatob2024model}
Elvis Dohmatob, Yunzhen Feng, and Julia Kempe.
\newblock A tale of tails: Model collapse as a change of scaling laws.
\newblock In \emph{ICML}, 2024.

\bibitem[Feng et~al.(2024)]{feng2024factuality}
Shangbin Feng, Vidhisha Balachandran, Yuyang Bai, and Yulia Tsvetkov.
\newblock FactScore: Fine-grained atomic evaluation of factual precision in long form text generation.
\newblock In \emph{EMNLP}, 2024.

\bibitem[Groeneveld et~al.(2024)]{groeneveld2024olmo}
Dirk Groeneveld, Iz Beltagy, Pete Walsh, Akshita Bhagia, Rodney Kinney, Oyvind Tafjord, Ananya Harsh Jha, Hamish Ivison, Ian Magnusson, Yizhong Wang, et~al.
\newblock {OLMo}: Accelerating the science of language models.
\newblock In \emph{ACL}, 2024.

\bibitem[Hans et~al.(2024)]{hans2024binoculars}
Abhimanyu Hans, Avi Schwarzschild, Valeriia Cherepanova, Hamid Kazemi, Aniruddha Saha, Micah Goldblum, Jonas Geiping, and Tom Goldstein.
\newblock Spotting {LLMs} with Binoculars: Zero-shot detection of machine-generated text.
\newblock In \emph{ICML}, 2024.

\bibitem[Kirchenbauer et~al.(2023)]{kirchenbauer2023watermark}
John Kirchenbauer, Jonas Geiping, Yuxin Wen, Jonathan Katz, Ian Miers, and Tom Goldstein.
\newblock A watermark for large language models.
\newblock In \emph{ICML}, 2023.

\bibitem[Mitchell et~al.(2023)]{mitchell2023detectgpt}
Eric Mitchell, Yoonho Lee, Alexander Khazatsky, Christopher~D Manning, and Chelsea Finn.
\newblock {DetectGPT}: Zero-shot machine-generated text detection using probability curvature.
\newblock In \emph{ICML}, 2023.

\bibitem[Qwen Team(2025)]{qwen2025qwen25}
Qwen Team.
\newblock Qwen2.5: A party of foundation models.
\newblock \emph{arXiv preprint arXiv:2412.15115}, 2025.

\bibitem[Raffel et~al.(2020)]{raffel2020c4}
Colin Raffel, Noam Shazeer, Adam Roberts, Katherine Lee, Sharan Narang, Michael Matena, Yanqi Zhou, Wei Li, and Peter~J Liu.
\newblock Exploring the limits of transfer learning with a unified text-to-text transformer.
\newblock \emph{JMLR}, 21(140):1--67, 2020.

\bibitem[Shumailov et~al.(2024)]{shumailov2024model}
Ilia Shumailov, Zakhar Shumaylov, Yiren Zhao, Nicolas Papernot, Ross Anderson, and Yarin Gal.
\newblock {AI} models collapse when trained on recursively generated data.
\newblock \emph{Nature}, 631:755--759, 2024.

\bibitem[Tulchinskii et~al.(2024)]{tulchinskii2024diveye}
Eduard Tulchinskii, Kristian Kuznetsov, Laida Kushnareva, Daniil Cherniavskii, Sergey Nikolenko, Evgeny Burnaev, Serguei Barannikov, and Irina Piontkovskaya.
\newblock Intrinsic dimension estimation for robust detection of {AI}-generated texts.
\newblock In \emph{NeurIPS}, 2024.

\bibitem[Wettig et~al.(2024)]{wettig2024qurating}
Alexander Wettig, Aatmik Gupta, Saumya Malik, and Danqi Chen.
\newblock {QuRating}: Selecting high-quality data for training language models.
\newblock In \emph{ICML}, 2024.

\bibitem[Zhang et~al.(2024)]{zhang2024minkpp}
Yijian Zhang, Jiaao Chen, Diyi Yang, and Thomas Hartvigsen.
\newblock Min-{K}\%++: Improved baseline for detecting pre-training data from large language models.
\newblock In \emph{EMNLP}, 2024.

\end{thebibliography}

\end{document}
